\fsection{Ej 4 Pág 142}
\fsubsection
[\textcolor{waveBlue2}{Tarea: Copias de Seguridad}]
{\textcolor{waveBlue2}{\boldit{Tarea: Copias de Seguridad}}}
{
	Reflexionad sobre los datos almacenados en vuestros equipos y dispositivos que no os gustaría
	perder (fotografías, vídeos, documentos\ldots). ¿En qué tipo de soportes los almacenáis?
	¿Qué supondría perderlos? ¿Lleváis a cabo alguna estrategia de copias de seguridad? ¿Utilizáis
	la regla del 3-2-1? ¿Creéis que es importante disponer de los datos duplicados en diferentes
	soportes?
}

\begin{solution}

	\textbf{1. Datos almacenados y su valor}

	En los dispositivos personales se acumulan datos de difícil o imposible recuperación:
	fotografías y vídeos personales o familiares, documentos académicos y trabajos de curso,
	configuraciones personalizadas de aplicaciones, contraseñas almacenadas, proyectos propios
	y archivos de trabajo. La pérdida de estos datos por un fallo del sistema, un robo, un
	incendio o un ataque de ransomware podría suponer un impacto irreversible tanto sentimental
	como académico o profesional.

	\medskip

	\textbf{2. Soportes habituales de almacenamiento}

	\begin{itemize}
		\item \textbf{Almacenamiento local:} disco duro interno del equipo (HDD/SSD). Es el más
		      común, pero también el más vulnerable a fallos físicos, robos o desastres.
		\item \textbf{Almacenamiento externo:} discos duros externos y pendrives. Ofrecen una copia
		      adicional, aunque también son susceptibles de pérdida física o fallo mecánico.
		\item \textbf{Almacenamiento en la nube:} servicios como Google Drive, OneDrive o Dropbox.
		      Protegen frente a fallos del hardware local, pero dependen de conexión a internet y de
		      terceros.
	\end{itemize}

	\medskip

	\textbf{3. Consecuencias de perder los datos}

	Un fallo del sistema sin copia de seguridad puede suponer: pérdida definitiva de recuerdos
	personales, necesidad de repetir trabajos académicos desde cero, pérdida de proyectos en
	desarrollo, y en un contexto profesional, incumplimiento de obligaciones legales de
	conservación de datos (RGPD, normativa mercantil y fiscal).

	\medskip

	\textbf{4. La regla 3-2-1}

	La regla 3-2-1 es la estrategia de backup más recomendada por expertos en seguridad:

	\begin{itemize}
		\item \textbf{3} copias de los datos (el original más dos copias de seguridad).
		\item Almacenadas en \textbf{2} tipos de soporte diferentes (p.~ej., disco interno y disco
		      externo, o disco interno y nube).
		\item Con al menos \textbf{1} copia en una ubicación física diferente o fuera de línea
		      (offsite), para protegerse frente a robos, incendios o inundaciones que afecten a un
		      mismo lugar.
	\end{itemize}

	\medskip

	\textbf{5. Valoración personal}

	Sí, considero fundamental aplicar una estrategia de copias de seguridad. Un único soporte
	nunca ofrece garantías suficientes: los discos duros tienen una vida útil limitada, los
	pendrives se extravían y los servicios en la nube pueden sufrir cortes o cambios de política.
	La regla 3-2-1 representa un equilibrio razonable entre coste y protección. En el ámbito
	profesional, no aplicarla puede constituir además un incumplimiento del deber de diligencia
	en la custodia de datos personales bajo el RGPD.

\end{solution}
